\section{Difesa e Validazione di
BlueAgent}\label{difesa-e-validazione-di-blueagent}

\subsection{Scelte di Piattaforma e
Infrastruttura}\label{scelte-di-piattaforma-e-infrastruttura}

\textbf{Perché Debian 13?} Offre un kernel moderno e stabile, garantendo
un'ottima compatibilita tra i vari pacchetti e le componenti utilizzate.

\textbf{Perche' ZFS?}: e' un file system generalmente ritenuto stabile e
affidabile che offre un sistema di CoW (Copy-on-Write) a basso costo. E'
stato scelto rispetto ad alternative come BTRFS poiche, ZFS e' ritenuto
più maturo e stabile.

\textbf{Perche' Incus?} A differenza di Docker, offre i container di
sistema, ossia container con un init system completo (e quindi con un
proprio processo \passthrough{\lstinline!init!} con
\passthrough{\lstinline!PID!} 0). Permette ai concorrenti di operare
realisticamente, accedendo ai servizi di sistema ed utilizzare utility
complicate come \passthrough{\lstinline!apt!}, seppur in ambiente
containerizzato. Inoltre Incus supporta nativamente l'uso di ZFS e di
OVN.

\subsection{Difesa dell'Architettura}\label{difesa-dellarchitettura}

Il vantaggio piu grande che porta BlueAgent e' l'invisibilita totale dei
giocatori rispetto all'Oracolo, ribadiamo dunque che e' impossibile per
i partecipanti intercettare ed ascoltare traffico di rete ed accedere ai
container tramite ssh a causa di come sono impostate le tabelle di
routing nell'host. Inoltre usare il meccanismo CoW riduce notevoilmente
il tempo iniziale richiesto per eseguire il setup dell'infrastruttura di
rete mentre l'uso di OVN facilita la configurazione di rete.

\subsection{Critica e Limiti
Architetturali}\label{critica-e-limiti-architetturali}

Nonostante i vantaggi, l'architettura presenta dei compromessi.

\begin{itemize}
\tightlist
\item
  \textbf{MTU Blackhole per UDP:} il traffico tra giocatore e container
  attraversa il tunnel WireGuard, con \(\text{MTU} 1420\).Mentre il
  problema viene risolto per TCP tramite l'MSS clamping, i pacchetti UDP
  con flag \emph{Don't Fragment} potrebbero essere scartati dal kernel
  in caso i pacchetti assumano una MTU pari a 1500 come in Ethernet.
\item
  \textbf{Limite scalabilità a 254 team:} L'indirizzamento
  deterministico \passthrough{\lstinline!10.k.x.x!} vincola \(k\) a un
  singolo byte (\(0 \le k \le 254\)), quindi ad un massimo di
  \(\text{255}\) team.
\item
  \textbf{Concorrenza e blocco multi-nodo:} Il meccanismo
  \passthrough{\lstinline!TwoTierLock!}
  (\passthrough{\lstinline!asyncio.Lock!} combinato con
  \passthrough{\lstinline!flock!}) è confinato al singolo nodo. Per
  scalare l'architettura su più nodi bisogna introdurre l'uso di un
  gestore di lock distribuito come Redis.
\item
  \textbf{Overhead in fase di Verifica}: gestire i cloni raddopia
  potenzialmente il consumo di CPU e di RAM per un singolo team. Al caso
  peggiore tutti i team potrebbero richiedere contemporaneamente la
  verifica dei loro container penalizzando le prestazioni dell'host.
\end{itemize}
